\documentclass[11pt]{article}
\usepackage[utf8]{inputenc}
\usepackage{amsmath,amssymb}
\usepackage{authblk}
\usepackage[margin=1in]{geometry}
\usepackage[hidelinks]{hyperref}
\usepackage{xurl}
\usepackage{setspace}

\title{Magnetism without the Complex Plane}
\author[1]{Claes Johnson}
\affil[1]{KTH Royal Institute of Technology, SE-100 44 Stockholm, Sweden \\
\texttt{claesjohnson@gmail.com}}
\date{\today}

\begin{document}
\maketitle

\begin{abstract}
Orbital magnetism is routinely presented as intrinsically complex: a circulating current needs a phase, and a
phase means a complex wave function $\psi=\sqrt\rho\,e^{iS}$. We separate, within RealQM, what is genuinely
quantum here from what is mere bookkeeping. From the charge density \emph{alone} --- a single real field --- no
magnetism can come, since the current $\mathbf j=\operatorname{Im}(\psi^*\nabla\psi)$ vanishes identically for
real $\psi$: a lone amplitude carries no motion. But complex numbers are not the only way to carry motion.
Continuum mechanics never uses them, yet always carries \emph{two} real fields, a density $\rho$ and a velocity
$\mathbf v$; carried the same way here, the whole of orbital magnetism is real. The current is
$\mathbf j=\rho\,\mathbf v$, a real circulating flow; it costs a real kinetic energy $\int\tfrac12\rho|\mathbf
v|^2$; it carries a real angular momentum and the ordinary magnetostatic moment; and the magnetic state is
obtained by minimising the static energy augmented by that flow energy at fixed angular momentum. No $i$ enters
any of it: magnetism is a real charge vortex, the $(\rho,\mathbf v)$ of continuum mechanics an
\aa ngstr\"om across. The complex $\psi=\sqrt\rho\,e^{iS}$ is exactly these two real fields repackaged into one
symbol, the phase being the velocity potential. The single thing the real theory cannot supply on its own is
the \emph{integer}: single-valuedness of $e^{iS}$ forces the quantised circulation $\oint\mathbf v\cdot
d\boldsymbol\ell=2\pi m$, $m\in\mathbb Z$, which a purely real formulation must impose as a separate condition
(the Onsager--Feynman vortex quantisation of real superfluids). What is unavoidable there is the \emph{circle},
the periodicity of an angle, not the complex plane. Magnetism's dynamics, energy and moment are real; only the
quantum integer needs the periodic phase.
\end{abstract}

\setstretch{1.06}

\section{The question}

Magnetism is charge in motion, and motion in quantum mechanics is carried by a phase: the standard route to an
orbital magnetic moment runs through a complex, phase-winding wave function $\psi=R(\mathbf x)\,e^{im\varphi}$,
whose current and moment follow from the imaginary part of $\psi^*\nabla\psi$. It is natural to conclude that
the complex wave function is \emph{necessary} for magnetism --- that here, at least, RealQM must leave the real,
static picture behind and become complex.

The conclusion is half right and worth taking apart, because the halves belong to different things. Something
beyond the bare density is indeed needed. But that something is a \emph{velocity}, not a complex number; and
complex arithmetic, where it appears, is a convenient packaging of two real fields, not an irreducible
ingredient. We make the separation explicit: what magnetism needs, what the complex form merely bundles, and
the one genuinely quantum residue that the real theory cannot generate on its own.

\section{From the density alone: nothing}

Let the only field be the real amplitude, $\psi=R$, so that the charge density is $\rho=q R^2$. The
quantum-mechanical current is
\[
\mathbf j=\frac{q}{m_e}\,\operatorname{Im}(\psi^*\nabla\psi),
\]
and for real $\psi$ this is \emph{identically zero}. A single real field has no current in it and therefore no
magnetic moment. This is the true content of the slogan ``magnetism needs the complex form'': from the density
by itself, nothing circulates, because a lone amplitude encodes no motion. Whatever carries magnetism must be a
\emph{second} field, over and above $\rho$.

\section{Two real fields: $(\rho,\mathbf v)$ continuum mechanics}

Complex numbers are one way to carry a second field, but not the only way. Macroscopic continuum mechanics
carries motion with no complex numbers at all: it works throughout with a density $\rho$ and a velocity
$\mathbf v$. Carry the same pair in RealQM and the whole of orbital magnetism is real. The current is a real
circulating flow,
\begin{equation}
\mathbf j=\rho\,\mathbf v,\qquad \nabla\!\cdot\!\mathbf j=0\ \text{(stationary)},
\label{eq:j}
\end{equation}
a charge vortex. The flow costs a real kinetic energy,
\begin{equation}
T_{\text{flow}}=\int \tfrac12\,\rho\,|\mathbf v|^2\,d^3x=\int \frac{|\mathbf j|^2}{2\rho}\,d^3x ,
\label{eq:T}
\end{equation}
it carries a real angular momentum and the ordinary magnetostatic moment of a current distribution,
\begin{equation}
L_z=\int \rho\,(\mathbf x\times\mathbf v)_z\,d^3x ,\qquad
\boldsymbol\mu=\tfrac12\int \mathbf x\times\mathbf j\,d^3x ,
\label{eq:muL}
\end{equation}
and the magnetic state is the stationary point of the \emph{static} RealQM energy augmented by the flow energy
at fixed angular momentum,
\begin{equation}
\min\ \big\{\,E_{\text{static}}[\rho]+T_{\text{flow}}[\rho,\mathbf v]\ :\ L_z\ \text{fixed}\,\big\},
\qquad\text{or, in a field,}\quad \min\ \big(E-\boldsymbol\mu\!\cdot\!\mathbf B\big).
\label{eq:var}
\end{equation}
No $i$ enters \eqref{eq:j}--\eqref{eq:var}. This is a real charge vortex, precisely the $(\rho,\mathbf v)$ of
continuum mechanics scaled down to atomic size --- the same object a fluid dynamicist would write for a
swirling incompressible flow, now made of charge and confined to an atom.

\section{What the complex form packages}

Where, then, does the familiar complex wave function come in? Write it in polar (Madelung) form,
\[
\psi=\sqrt{\rho}\;e^{iS},\qquad \mathbf v=\frac{\nabla S}{m_e},
\]
and the identification is immediate: the modulus is the (root) density, and the phase $S$ is nothing but the
\emph{velocity potential}, its gradient the velocity. The current $\mathbf j=\rho\,\mathbf v$ of \eqref{eq:j}
is exactly $\operatorname{Im}(\psi^*\nabla\psi)$ times the charge-to-mass factor. So $\psi=\sqrt\rho\,e^{iS}$
is not a new ingredient at all: it is the two real fields $(\rho,\mathbf v)$ of \S3 bundled into a single
complex symbol. The ``$i$'' is bookkeeping --- a compact way to carry an amplitude and a flow together --- and
the physics is the flow. (This is also why, numerically, the honest integrator uses the Cartesian real pair
$(u,v)=(\operatorname{Re}\psi,\operatorname{Im}\psi)$, smooth through nodes, rather than the polar form, which
is singular exactly on the vortex axis where the amplitude vanishes.)

\section{The one quantum residue: the integer}

There is a single thing the real $(\rho,\mathbf v)$ theory cannot supply on its own, and it is worth naming
precisely. A winding state has phase $S=m\varphi$ about an axis, giving in cylindrical coordinates
$(s,\varphi,z)$ the azimuthal velocity $\mathbf v=(m/m_e s)\,\hat{\boldsymbol\varphi}$ and hence
$\mu_z=-m\mu_B$, quantised in $m$. Why must $m$ be an \emph{integer}? Because $S$ is an \emph{angle}: the phase
is defined only modulo $2\pi$, equivalently $e^{iS}$ must be single-valued, and this forces the circulation to
be quantised,
\begin{equation}
\oint \mathbf v\cdot d\boldsymbol\ell=\frac{2\pi m}{m_e},\qquad m\in\mathbb Z ,
\label{eq:circ}
\end{equation}
and with it $L_z=m\hbar$ and the moment in whole units of $\mu_B$. A purely real formulation carrying only
$(\rho,\mathbf v)$ would have every ingredient of the vortex except this: it would have to \emph{impose}
\eqref{eq:circ} as a separate quantisation condition --- which is exactly what one does for the quantised
vortices of a real superfluid (Onsager and Feynman). And what is genuinely required there is the
\emph{circle}, the periodicity of an angle, not the complex plane as such. The density must also vanish on the
axis (a nodal line, the hollow vortex core), since $\mathbf v\sim 1/s$ would otherwise diverge; that too is a
consequence of the winding, not of complex arithmetic.

\section{Consequences}

The upshot sharpens RealQM's guiding slogan --- the same 3D form as macroscopic continuum mechanics, only
smaller. Continuum mechanics is $(\rho,\mathbf v)$; a steady magnetic current is therefore \emph{native} to the
real theory, not an import from a complex one. Orbital magnetism --- moments, angular momentum, the Zeeman
response to an applied field, diamagnetic counter-currents --- is real continuum mechanics; the sole quantum
input is the integer winding number, and that input is the periodicity of an angle.

Two boundaries keep the account honest. First, this covers the \emph{steady} member of the ``complex'' family
only. \emph{Radiation} is the \emph{oscillating} member: a transition makes the density itself oscillate in
time at $E_n-E_m$, and a genuinely time-varying density does require the complex, time-dependent form, since a
real currentless density is static and dark~\cite{Spectra}. The temporal phase there is real physics (imposed
by the field), whereas the spatial phase of magnetism is, as shown, repackaged real flow. Second, single-%
electron \emph{spin} ($g=2$, Stern--Gerlach) is a further residue, reached from a spinor structure and not
from the scalar $(\rho,\mathbf v)$ picture of orbital magnetism~\cite{Magnetism}. Within those boundaries the
statement stands: magnetism's dynamics, energy and moment are real; only the quantum integer needs the periodic
phase.

\section{What is new, what is not, and the boundary}

Because the phenomenology here is standard, it is worth stating plainly what RealQM does and does not add, so the
contribution is not mistaken either way.

\paragraph{What is genuinely new is the ontology.} In the usual account the orbital moment is an operator
$\hat{\mathbf L}$ on an amplitude, and the ``current'' is a \emph{probability} current --- physical for one
particle, and otherwise a construct on $3N$-dimensional configuration space. RealQM makes magnetism a literal
\emph{charge fluid} $(\rho,\mathbf v)$ circulating in ordinary three-dimensional space, one genuine current per
electron for any $N$: a real charge vortex, the same object continuum mechanics writes for a swirling flow, an
\aa ngstr\"om across~\cite{RealQM}. From this follow two clarifications with no counterpart in the standard
telling. First, the complex wave function is, for \emph{steady} magnetism, mere bookkeeping: what is needed is a
velocity, not an $i$, and $\psi=\sqrt\rho\,e^{iS}$ merely bundles $(\rho,\mathbf v)$, so the only irreducibly
quantum residue is the integer $m$ --- forced by the \emph{circle} (the periodicity of an angle), not by the
complex plane. Second, this locates magnetism precisely against radiation: both belong to the ``complex family,''
but magnetism is a \emph{spatial} winding (steady, real flow in disguise) whereas radiation is a \emph{temporal}
oscillation that genuinely requires the complex time-dependent form, a real currentless density being
dark~\cite{Spectra}. Same family, opposite time-status.

\paragraph{What is not new is the phenomenology.} RealQM \emph{reproduces}, it does not rewrite, the standard
results: the quantised circulation, the gyromagnetic relation $\boldsymbol\mu=(q/2m_e)\mathbf L$, the Zeeman
splitting, and Langevin diamagnetism $\chi\propto-\langle r^2\rangle$. The hydrodynamic $(\rho,\mathbf v)$ form
is Madelung's~\cite{Madelung}; quantised vortices are Onsager's and Feynman's~\cite{Onsager,Feynman}. No new
magnetic number or law is claimed here; the novelty is in \emph{what these quantities are}, not in their values.

\paragraph{The boundary is collective magnetism.} What RealQM does \emph{not} supply is exchange-driven
magnetism --- the singlet--triplet splitting, ferromagnetic ordering --- which requires spin exchange that
geometric non-overlap does not deliver. Single-electron spin is reachable as a further residue ($g=2$, without
relativity), but the collective effects are not, and this note claims no perspective that resolves
them~\cite{Magnetism}. The account of orbital magnetism above is complete and real; the exchange sector remains
open, and is named as such.

\begin{thebibliography}{9}
\bibitem{RealQM} C.~Johnson, \emph{Quantum Mechanics as Multiphase 3D Continuum Mechanics}, manuscript,
2026; \url{https://claes542.github.io/RealMolecule/gallery.html}.
\bibitem{Magnetism} C.~Johnson, \emph{Magnetism in RealQM: Currents, the Nuclear Electron, and the Spin
Residue}, manuscript, 2026.
\bibitem{ComplexRealQM} C.~Johnson, \emph{Complex Time-Dependent RealQM: Currents, Radiation, and Magnetism
from the Real-Space Continuum}, manuscript, 2026.
\bibitem{Spectra} C.~Johnson, \emph{RealQM Atomic Spectra: Radiation as Real Charge Oscillation}, manuscript,
2026.
\bibitem{Time} C.~Johnson, \emph{Real Time Dependence in RealQM}, manuscript, 2026.
\bibitem{Madelung} E.~Madelung, \emph{Quantentheorie in hydrodynamischer Form}, Z.\ Phys.\ \textbf{40}, 322
(1927).
\bibitem{Onsager} L.~Onsager, \emph{Statistical hydrodynamics}, Nuovo Cimento \textbf{6} (Suppl.), 279 (1949).
\bibitem{Feynman} R.~P.~Feynman, \emph{Application of quantum mechanics to liquid helium}, in \emph{Progress
in Low Temperature Physics}, Vol.~1, North-Holland, 1955.
\end{thebibliography}

\end{document}
